% Set spellcheck for TexStudio
%!TeX spellcheck = en_US
%!TeX root = ../thesis.tex

\chapter{Implementation} \label{Implementation}
Here we want to describe the developed algorithms. As already mentioned in the introduction the goal is to efficiently calculate the union $U$
where every processor $r$  provides a subset $S_r\subset \tilde{S}$ from a known set $\tilde{S}$ of fixed size. The union
\begin{align*}
	U = \bigcup_{r=0}^{p-1} S_r
\end{align*}
should be available for every processor. $p$ is here the number of processes. \\
A straight forward approach to tackle this problem is to create a bitmap for each process, give the elements of $\tilde{S}$ an explicit ordering
and map each element to a bit of the bitmap. An element $e$ is then defined as being in the subset $S_r$, if the the corresponding 
bit is set to 1 and set to 0 otherwise. The union can then easily be calculated by performing a bitwise OR-operation over all 
bitmaps from all processes as the union over all subset $S_r$ is defined as
\begin{align*}
	U = \{x \in C | x \in S_0 \lor x \in S_1 \lor \cdots \lor x \in S_{p-1} \}.
\end{align*}
This can be done using an Allreduce function. This bitmap needs to cover all elements of the original set $\tilde{S}$. Therefore it needs to take at least $N=|\tilde{S}|$ amount of data. As a word size of 64-bit integer is quite common on many architectures, we need $\ceil{\frac{N}{64}}$ 64-bit integer to represent a subset $S_r$. This can be wasteful, if the cardinality of these subsets $S_r$ is quite small and most bits of the bitmap are set to 0. We introduce here a simple sparse format to save memory and therefore the amount of data for communication. By using unsigned 32-bit integer, we can represent subsets of $\tilde{S}$ as an explicit list of the selected integer. Here we only can represent $2^{32}$ distinct elements so the cardinality of $\tilde{S}$ is bounded by $2^{32} \approx 4.2\cdot10^9$. The benefit is now that the size of the sparse representation only depends on the cardinality of the subset $S_r$. We can calculate the memory footprint $m_\mathrm{Sparse}$ of this representation as being 
\begin{align*}
	m_\mathrm{Sparse} = \varphi \cdot \qty{32}{\bit},
\end{align*}
where $\varphi = |S_r|$ is the cardinality of $S_r$. To calculate the set union $U$ one needs to concatenate all these distributed lists $S_r$ ($S_r$ is here represented as a list). This can be done by an Allgather. One needs to keep in mind that the result $U$ may contain elements twice. One can either decide that this not a problem or change representation and convert the list into a bitmap. A bitmap can not contain elements twice and therefore duplicates are removed. \\
We chose algorithms that scale well for a large number of processors $p$ and large number of elements of the base set $N$. For the Allreduce these properties are best fulfilled by any of the Reduce\_Scatter-Allgather algorithms. We moreover can identify that the communication pattern of the Reduce\_Scatter phase always resembles that of an Allgather. This can easily be seen for the Butterfly algorithm and the Ring algorithm. For the approximate halfing pattern there is also an Allgather known \cite{allgatherTraeffOld}.\todo{Add section for this Allgather implementation in Introduction chapter.} The problem here is that the skip sequence is used in reversed order compared to the one used in the Reduce\_Scatter phase from Section \ref{sec:approxHalfRedScat}. For this pattern is a new adaptation found which is shown in Section \ref{sec:irregularImpl}. \\
The general idea is therefore to overlay an Allgather using the sparse representation over the Reduce\_Scatter phase of the corresponding Allreduce algorithm. When it is more efficient to use the bitmap variant one switches to this representation and finishes the Allreduce algorithm. This is possible as the subsolutions of the Allgather span the same processes as the subsolutions of the Reduce\_Scatter do. As the sparse representation holds the complete information of a ranks input and the dense representation uses only blocks of it, one does not lose any information, if we switch from sparse to dense. This implies of course that a switch from the bitmap variant back to the sparse one is not permitted. How and why these ideas work is proved and explained in the upcoming respective sections.\\
If the sparse variant is still used for the last iteration of the Reduce\_Scatter phase the algorithm performed de facto an Allgather and therefore the second phase (Allgather phase of the Allreduce) is not necessary. Moreover the sparse variant may be shorter in message size than the blocks send in Reduce\_Scatter phase of the dense representation. So here it is also possible that the total amount of sent data can be reduced.\\
This is also relevant in choosing the foundation Allreduce algorithms for the newly developed ones. The Ring algorithm sends in each iteration only one block of data. The algorithms that are based on recursive (approximately) halfing send in iteration $i$ (roughly) $p/2^i$ blocks. Sending a sparse representation is therefore more efficient in the later algorithms compared to the Ring algorithm. We here replace the message size of multiple blocks instead of only one. This is the reason we chose recursive (approximately) halfing algorithms as base Allreduce algorithms over the Ring algorithm.\\
Moreover we chose two different approaches in communication schemes. One is chosen to be the Butterfly scheme as we may expect better performance of this scheme as it can be arranged in a way that possibly faster intra nodal communication can be exploited in some situations (see run time estimation in Section \ref{sec:DisseminationAlgo}). A big drawback of this scheme is of course that it is limited to only work correctly on a setup where $p$ is a power of two. So on the other hand we want to propose an adaption of the approximate halfing scheme to create an algorithm for arbitrary number of processes. These two adaptations are shown in the next sections.

\section{An Algorithm for $2^k$-Number of Processes}
As mentioned before the base communication scheme is the Butterfly scheme. In each iteration there is communication between pairs of processors, meaning, if rank $r$ sends a message to its partner $\hat{r}$, then $\hat{r}$ sends in this iteration also a message to $r$. The communication partner is found by flipping the $i$-th bit where $i$ is the iteration. The difference between IDs of communicating ranks is from now on called 'shift'. 'shift' also used as variable name in the according pseudocode \ref{algo:ExactHalfingSetUnion}. As this variable is set to $p/2$ in the first iteration ($i=0$) and is halved every iteration, we can identify its value as being $p/2^i$. Moreover we see in the implementation that the algorithm consist out of an adapted Reduce\_Scatter phase and an Allgather with the exact same communication scheme like Algorithm \ref{algo:redScatterAndAllreduce}. To be precise it is identical, if the own representation as well as the received representation is always dense. The receiving part for this situation is hidden behind the function $\Call{combineDenseDense}{ }$. This is done for all possible combinations to get a better overview of the control flow. The functionality of all cases is explained a little bit later in this section. Firstly we want to describe the control flow to get a better understanding of the general ideas.\\
We see that in the adapted Reduce\_Scatter phase there are two different possibilities which kind of data is send. If  variable '$isDense$' is true, process $r$ holds currently data in dense representation and therefore sends dense data blocks accordingly to the recursive halfing scheme. An add-on is here that the message is also marked with an label "dense". If  variable '$isDense$' is false, process $r$ holds currently data in sparse representation and sends its whole array of sparse elements. In this case the message gets obviously tagged with "sparse". After rank $r$ has posted a send it is ready to receive some data. This is done in one of the "combine" functions depending on the state of $S_r$ and the representation of the received message. 
\begin{algorithm}
	\caption{Distributed set union using recursive exact halfing Allreduce with sparse adaptation.\\
		$r$: rank Id of the process\\
		$p$: number of processes\\
		$I_s$: input array with sparse representation of $S_r$\\
		$sparseSize$: length of sparse list\\
		$I_d$: input array with dense representation of $S_r$\\
		$bitarraySize$: count of words for bitarray\\
		$isDense$: bool to indicate, which representation holds} \label{algo:ExactHalfingSetUnion}
	\begin{algorithmic}[1]
		\Ensure $p = 2^k$ for some $k$ \Comment{Only works for power of 2 number of ranks}
		\Procedure{SetUnion}{$r, p, I_s, sparseSize, I_d, bitarraySize, isDense$}
		\State critSize $\gets$ \Call{getCritSize}{$bitarraySize$}
		\Comment{Calculate Critical size}
		
		\For{(shift $\gets p/2$, shift $\geq1$, shift $\gets$ shift/2)} \Comment{Reducescatter phase}
			\State commNeigh $\gets r$ XOR shift \Comment{XOR is bitwise XOR}
			\If{$isDense$} \Comment{Send Data}
				\State startBlock $\gets$ (commNeigh/shift) $\cdot$ shift
				\State endBlock $\gets$ startBlock + shift
				\State Send($I_d$[startBlock, $\dots$, endBlock - 1], "dense") 
				\Statex\Comment{Send block. Mark message as "dense"}
			\Else
				\State Send($I_s$[0, $\dots$, $sizeSparse$-1], "sparse") 
				\Statex\Comment{Send whole sparse array. Mark message as sparse}
			\EndIf
			\State status $\gets$ \Call{Probe}{commNeigh} 
			\Statex\Comment{Get incomming message type and size information}
			\If{(status.tag = "dense")} 
			\Statex\Comment{Combine incomming message with own representation}
				\If{$isDense$}
				\State\Call{combineDenseDense}{}
				\Else 
				\State\Call{combineDenseSparse}{}
				\EndIf
			\ElsIf{(status.tag = "sparse")} 
				\If{($isDense$)}
				\State\Call{combineSparseDense}{}
				\Else
				\State\Call{combineSparseSparse}{}
				\EndIf
			\EndIf			
		\EndFor
		\If{($isDense$ = false)} \Comment{Return early, if variant is still sparse}
			\State \Return
		\EndIf
		\State\Call{Allgather}{$I_d$} \Comment{Allgather phase}
		\EndProcedure
	\end{algorithmic}
\end{algorithm}
\\We also want to mention here that these sends differ from the $\Call{Send}{ }||\Call{Recv}{ }$ procedure used in all algorithms before. The sending and receiving is now not done by the same function where everything is handled by an already implemented library. When called separately the send obviously needs to be called before the receive and usually is 
\emph{blocking}. Blocking means here the function only returns after the message is completely send. To complete a send another process needs to expect some data from the sending rank. This can lead to a deadlock situation. If process $r$ sends a message to process $\hat{r}$ and process $\hat{r}$ wants to do the same the other way around, both of them can not return from the $\Call{Send}{ }$ function as none of them is already ready to expect a message. So in this implementation we use non-blocking sends. For this some extra handling is needed. As the send function returns immediately one does not want to change the buffer array that needs to be send until the message is completely sent. To guarantee safe usage of the send array the function $\Call{Wait}{ }$ is used. It blocks further execution of the program until the corresponding buffer is safe to use again. In \gls{mpi} this actually handled by an extra request object, but we want to keep notation simple and define the input parameter of $\Call{Wait}{ }$ as the corresponding input array. \\
Looking at the input arrays $I_s$ and $I_d$ we want to give some requirements about them. $I_s$ is used to hold information of the subset $S_r$ rank $r$ has in the sparse format. It needs to be filled, if $isDense$ is set to false. It consists out of 32-bit integer and holds $sparseSize$ elements that correspond to the elements of $S_r$. The array needs have enough memory allocated to hold two times 'critSize' elements. When indexing this array we refer to single elements of this array. The same conditions holds for the extra buffer array $\Tilde{I}_s$. $I_d$ holds the bitarray representation of $S_r$ and needs to be filled, if $isDense$ is true. It is assumed to be empty (filled with only zeros), if the sparse variant is used as input. Its size is given by $bitarraySize$. The corresponding extra buffer $\Tilde{I}_d$ needs to be able to hold half of the biggest blocks of $I_d$. Both $I_d$ and $\Tilde{I}_d$ are assumed to be divided into $p$ blocks. These blocks are assumed to be roughly equally sized and hold approximately $bitarraySize/p$ amount of data. Indexing these arrays therefore refers to blocks of data. To assure these conditions the initialization of the arrays, calculation of the blocks and adding of elements to $S_r$ is hidden from the user in the original implementation.\\
Another function that needs to be explained is $\Call{Probe}{ }$. This function returns a status object for messages that rank $r$ receives from rank $\hat{r}$. This status object holds information about the tag that was added by the sending rank and how long the message is. The tag is needed to identify, if data received in a dense or a sparse variant. The size information is only needed, if a sparse array is received. In this case a processor does not know how many elements are added to the subset $S_{\hat{r}}$ it is receiving and therefore the message size is unknown. This problem does not arise in the dense variant as here are always a known amount of blocks of the array sent by $\hat{r}$ which can also be calculated by $r$ using $\hat{r}$ rank ID.\\
The next important feature we want to mention is the variable 'critSize'. It defines the boundary which decides, if $r$ should use the dense or sparse representation for its subset $S_r$. 'critSize' is compared against the size of $S_r$ $\varphi$. If $\varphi$ is bigger than 'critSize' process $r$ uses the dense variant and if it is smaller the sparse format is used. There can be different strategies on how to chose 'critSize' which is for now hidden behind $\Call{getCritSize}{ }$ and explained later.\\
We already explained which data is send in each iteration $i$. Now we want to focus how the received messages are processed in the different cases. As we see there are four different possibilities depending on the representation of the received message (identified by 'status.tag') and on the representation rank $r$ currently uses (identified by $isDense$). In all cases we somehow want to combine the (possible) different representation into a definitive one. We start with the case where the receiving message as well rank $r$s representation is dense. In Algorithm \ref{helper:combineDenseDense} we see the procedure for this case and can identify that it is nearly the same as the receive handling of the Reducescatter phase in the Butterfly recursive halfing Allreduce algorithm (compare to Algorithm \ref{algo:redScatterAndAllreduce}).
\begin{algorithm}
	\caption{– combine a received dense blocks into own dense array.\\
	Variable nomenclature is the same as in Algorithm \ref{algo:ExactHalfingSetUnion}.}
	\label{helper:combineDenseDense}
	\begin{algorithmic}[1]
		\Procedure{combineDenseDense}{$r$, shift, status}
			\State startBlockR $\gets$ ($r$/shift) $\cdot$ shift
			\State endBlockR $\gets$ startBlock + shift
			\State Recv($\tilde{I}_d$[0, $\dots$, endBlockR - startBlockR-1], "dense")\Comment{Receive in buffer}
			\State Wait($I_d$) \Comment{Wait until it is safe to use original array}
			\ForAll{i $\in$ \{startBlockR, $\dots$, endBlockR - 1\}}
			\State $I_d$[i] $\gets$ $I_d$[i] $\lor$  $\tilde{I}_d$[i - startBlock] \Comment{$\lor$ is here bitwise OR}
			\EndFor
		\EndProcedure
	\end{algorithmic}
\end{algorithm}
One small difference is here that the send and receive is not handled in the same function. The sending rank $\hat{r}$ used a non-blocking send and maybe also rank $r$ did and used $I_d$ in the same iteration for sending a message. Therefore we need to check, if buffer $I_d$ is already safe to use and wait, if this is not the case. The other difference is the reduce operation. It is now not a general associative operator $\oplus$ anymore, but instead the explicitly defined bitwise OR ($\lor$). This is of corse done as the bitwise OR performs a logically the set union of two bitmaps. Calculations of the block indices for the receiving block range are still exactly the same as in Algorithm \ref{algo:redScatterAndAllreduce}. Also the received blocks are "added" to the same range of blocks of $I_d$. We can summarize that after an iteration where $\Call{combineDenseDense}{ }$ is used a standard step of the Butterfly Reducescatter phase is performed.\\
In the next scenario rank $r$ receives again dense data blocks, but holds itself data in sparse representation. "status.tag" is therefore "dense", $isDense$ is false and $\Call{combineDenseSparse}{ }$ is executed. The corresponding procedure is shown in Algorithm \ref{helper:combineDenseSparse}. We see that the dense blocks are put at the dense array $I_d$ at positions [startBlockR, $dots$, endblockR-1]. The block indices are calculated the same way as they where for the purely dense case. Actually this part can be seen as equivalent to the procedure of putting the receiving blocks in a buffer and performing an logic OR to an empty dense array $I_d$. After the blocks are build into $I_d$ the elements of rank $r$s sparse 
need to be build in. This is done by iterating over all elements of the sparse array and set the corresponding bits of $I_d$ to 1 using $\Call{Bitmap_Set}{ }$.
\begin{algorithm}
	\caption{– combine a received sparse list into own dense array.\\
		Variable nomenclature is the same as in Algorithm \ref{algo:ExactHalfingSetUnion}.}
	\label{helper:combineDenseSparse}
	\begin{algorithmic}[1]
		\Procedure{Procedure}{$r$, shift, status, $isDense$}
			\State startBlockR $\gets$ ($r$/shift) $\cdot$ shift
			\State endBlockR $\gets$ startBlock + shift
			\State Recv($I_d$[startBlockR, $\dots$, endBlockR - 1], "dense")
			\ForAll{element $\in I_s$} \Comment{Add all old sparse elements in new dense array}
			\State \Call{Bitmap\_Set}{$I_d$, element}
			\EndFor	
			\State isDense $\gets$ true
		\EndProcedure
	\end{algorithmic}
\end{algorithm}



\begin{algorithm}
	\caption{Distributed set union using recursive exact halfing Allreduce with sparse adaptation.\\
		$r$: rank Id of the process\\
		$p$: number of processes\\
		$I_s$: input array with sparse representation of $S_r$\\
		$sparseSize$: length of sparse list\\
		$I_d$: input array with dense representation of $S_r$\\
		$bitarraySize$: count of words for bitarray\\
		$isDense$: bool to indicate, which representation holds} \label{algo:ExactHalfingAllreduce123}
	\begin{algorithmic}[1]
		\Ensure $p = 2^k$ for some $k$ \Comment{Only works for power of 2 number of ranks}
		\Procedure{SetUnion}{$r, p, I_s, sparseSize, I_d, bitarraySize, isDense$}
		\State critSize $\gets$ \Call{getCritSize}{$bitarraySize$}
		\Comment{Calculate Critical size}
		
		\For{(shift $\gets p/2$, shift $\geq1$, shift $\gets$ shift/2)} \Comment{Reducescatter phase}
		\State commNeigh $\gets r$ XOR shift \Comment{XOR is bitwise XOR}
		\If{isDense}
		\State startBlock $\gets$ (commNeigh/shift) $\cdot$ shift
		\State endBlock $\gets$ startBlock + shift
		\State Send($I_d$[startBlock, $\dots$, endBlock - 1], "dense") \Comment{Send block. Mark message as "dense"}
		\Else
		\State Send($I_s$[0, $\dots$, $sizeSparse$-1], "sparse") \Comment{Send whole sparse array. Mark message as sparse}
		\EndIf
		\\
		\State status $\gets$ \Call{probe}{commNeigh} \Comment{Get message type and size information}
		\If{(status.tag = "dense")} \Comment{Dense block is received}
		\State startBlock $\gets$ ($r$/shift) $\cdot$ shift
		\State endBlock $\gets$ startBlock + shift
		\If{$isDense$} \Comment{Own structure is also "dense"}
		\State Recv($\tilde{I}_d$[0, $\dots$, status.count - 1], "dense")\Comment{Receive in buffer}
		\State Wait($I_d$) \Comment{Wait until it is safe to use original array}
		\ForAll{i = startBlock, $\dots$, endBlock - 1}
		\State $I_d$[i] $\gets$ $I_d$[i] $\lor$  $\tilde{I}_d$[i - startBlock] \Comment{$\lor$ is here bitwise OR}
		\EndFor
		\Else \Comment{Own structure is "sparse"}
		\State Recv($I_d$[startBlock, $\dots$, startBlock + status.count - 1], "dense")
		\ForAll{element $\in I_s$} \Comment{Add all old sparse elements in new dense array}
		\State \Call{Bitmap\_Set}{$I_d$, element}		
		\EndFor					
		\State isDense $\gets$ true
		\EndIf
		\\
		\ElsIf{(status.tag = "sparse")} \Comment{Sparse array is received}
		\If{($isDense$)}
		\State Recv($\tilde{I}_s$ [0, $\dots$, ], "sparse")
		\ForAll{element $\in \tilde{I}_s$} \Comment{Add all new sparse elements in old dense array}
		\State \Call{Bitmap\_Set}{$I_d$, element}		
		\EndFor
		\Else
		\State Recv($I_s$[$sizeSparse$, $\dots$, $sizeSparse$+status.count-1], "sparse") \\ \Comment{Receive directly in sparse array behind the old elements}
		\State $sizeSparse \gets sizeSparse$ + status.count
		\If{($sizeSparse > critSize$)} \Comment{Transform into dense, if too big}
		\ForAll{element $\in I_s$} 
		\State \Call{Bitmap\_Set}{$I_d$, element}		
		\EndFor
		\EndIf					
		\EndIf
		\EndIf			
		\EndFor
		\\
		\If{($isDense$ = false)} \Comment{Return early, if variant is still sparse}
		\State \Return
		\EndIf
		\\
		\Call{Allgather}{$I_d$}
		\EndProcedure
	\end{algorithmic}
\end{algorithm}



\subsection{Concept and Explicit Implementation}
\subsection{Proof of Correctness}
\subsection{Runtime Estimation}


\section{An Algorithm for Arbitrary Number of Processes} \label{sec:irregularImpl}
\subsection{Concept and Explicit Implementation}
\subsection{Proof of Correctness}
\subsection{Runtime Estimation}

\section{A Naive Approach}


